% Part: first-order-logic
% Chapter: models-theories
% Section: theories

\documentclass[../../../include/open-logic-section]{subfiles}

\begin{document}

\olfileid{fol}{mat}{the}

\olsection{أمثلة لنظريات الرتبة الأولى}

\begin{ex}
تُعطى بديهيات نظرية الترتيبات الخطية الصارمة في اللغة~$\Lang L_<$
بالمجموعة
\begin{align*}
\{\quad & \lforall[x][\lnot x < x], \\
& \lforall[x][\lforall[y][((x < y \lor y <
    x) \lor x = y)]], \\
& \lforall[x][\lforall[y][\lforall[z][((x < y
      \land y < z) \lif x < z)]]] \quad \}
\end{align*}
وهي تجسد ال!!{structure}s المقصودة تجسيدًا كاملًا: فكل ترتيب خطي صارم
نموذج لنسق البديهيات هذا؛ وفي الاتجاه العكسي، إذا كانت $R$ ترتيبًا
خطيًا صارمًا على مجموعة $X$، فإن البنية $\Struct M$ التي فيها
$\Domain M = X$ و$\Assign{<}{M} = R$ نموذج لهذه النظرية.
\end{ex}

\begin{ex}
تُعطى بديهيات نظرية الزمر في لغة تحتوي $\Obj 1$ (!!{constant}) و$\cdot$
(!!{function} ثنائية المحل) بالآتي:
\begin{align*}
& \lforall[x][\eq[(x \cdot \Obj 1)][x]]\\
& \lforall[x][\lforall[y][\lforall[z][\eq[(x \cdot (y \cdot z))][((x
          \cdot y) \cdot z)]]]]\\
& \lforall[x][\lexists[y][\eq[(x \cdot y)][\Obj 1]]]
\end{align*}
\end{ex}

\begin{ex}
تُعطى بديهيات نظرية حساب بيانو بالجمل الآتية في لغة
الحساب~$\Lang L_A$:
\begin{align*}
& \lforall[x][\lforall[y][(\eq[x'][y'] \lif \eq[x][y])]]\\
& \lforall[x][\eq/[\Obj 0][x']]\\
& \lforall[x][\eq[(x + \Obj 0)][x]]\\
& \lforall[x][\lforall[y][\eq[(x + y')][(x + y)']]]\\
& \lforall[x][\eq[(x \times \Obj 0)][\Obj 0]]\\
& \lforall[x][\lforall[y][\eq[(x \times y')][((x \times y) + x)]]]\\
& \lforall[x][\lforall[y][(x < y \liff \lexists[z][\eq[(z' + x)][y])]]]\\
\intertext{ويضاف إليها جميع الجمل التي صورتها}
& (!A(\Obj 0) \land \lforall[x][(!A(x) \lif !A(x'))]) \lif \lforall[x][!A(x)]
\end{align*}
ولما كان عدد الجمل ذات الصورة الأخيرة لا نهائيًا، كان نسق البديهيات
هذا لا نهائيًا. وتسمى الصورة الأخيرة \emph{مخطط الاستقراء}. (والواقع
أن مخطط الاستقراء أعقد قليلًا مما أوحينا به هنا.)

والبديهية الأخيرة \emph{تعريف صريح} لـ~$<$.
\end{ex}

\begin{ex}
تؤدي نظرية المجموعات الصرفة دورًا مهمًا في أسس الرياضيات (وفي
فلسفتها). وتكون المجموعة صرفة إذا كانت جميع !!{element}sها مجموعات
صرفة أيضًا. ولذلك تعد المجموعة الخالية صرفة، أما المجموعة التي يكون
فيها شيء بوصفه !!a{element} وليس ذلك الشيء مجموعة فلا تكون صرفة. فالمجموعات الصرفة
هي إذن التي لا تتكون إلا من المجموعة الخالية، ولا يدخل في تكوينها أي
«عنصر أولي»، أي أي كائن ليس هو نفسه مجموعة.

ويمكن عد الآتي نسق بديهيات لنظرية المجموعات الصرفة:
\begin{align*}
& \lexists[x][\lnot \lexists[y][y \in x]]\\
& \lforall[x][\lforall[y][(\lforall[z][(z \in x \liff z \in y)] \lif
\eq[x][y])]]\\
& \lforall[x][\lforall[y][\lexists[z][\lforall[u][(u \in z \liff
(\eq[u][x] \lor \eq[u][y]))]]]]\\
& \lforall[x][\lexists[y][\lforall[z][(z \in y \liff \lexists[u][(z \in
        u \land u \in x)])]]]\\
\intertext{ويضاف إليها جميع الجمل التي صورتها} &
\lexists[x][\lforall[y][(y \in x \liff !A(y))]]
\end{align*}
تقول البديهية الأولى إن ثمة مجموعة لا !!{element}s لها (أي إن
$\emptyset$ موجودة)؛ وتقول الثانية إن المجموعات امتدادية؛ وتقول
الثالثة إن المجموعة $\{X, Y\}$ موجودة لأي مجموعتين $X$ و$Y$؛ وتقول
الرابعة إن المجموعة $\cup X$ موجودة لأي مجموعة $X$، حيث $\cup X$ هي
اتحاد جميع عناصر $X$.

وتسمى ال!!{sentence}s المذكورة أخيرًا، مجتمعة، \emph{مخطط الاستيعاب
الساذج}. وهي تقول، في جوهرها، إن المجموعة $\Setabs{x}{!A(x)}$ موجودة
لكل $!A(x)$؛ فتبدو لأول وهلة بديهية صادقة ومفيدة، بل ربما ضرورية.
وتسمى «ساذجة» لأنه يتبين أنها تجعل هذه النظرية غير قابلة للإرضاء: فإذا
أخذنا $!A(y)$ لتكون $\lnot y \in y$ حصلنا على ال!!{sentence}
\[
\lexists[x][\lforall[y][(y \in x \liff \lnot y \in y)]]
\]
وهذه ال!!{sentence} لا ترضى في أي !!{structure}.
\end{ex}

\begin{ex}
تُعد علاقة \emph{الجزء} علاقة أساسية في مبحث \emph{الجزء والكل}.
وكما توجد نظريات للمجموعات، توجد نظريات لعلاقة الجزء تصوغ بديهيًا
تصورات مختلفة (ومتعارضة أحيانًا) لهذه العلاقة.

تحتوي لغة مبحث الجزء والكل رمز محمول ثنائي المحل واحدًا~$\Obj P$،
و«تعني» $\Atom{\Obj P}{x, y}$ أن $x$ جزء من~$y$. وإذا وضعنا هذا
التأويل في اعتبارنا، سميت !!a{structure} لهذه اللغة \emph{بنية
للجزء}. ومن البدهي أن كل بنية لمحمول واحد ثنائي المحل لا تستحق هذه
التسمية حقًا. فلكي يكون ثمة احتمال لأن تجسد $\Assign{\Obj P}{M}$
«علاقة الجزء»، يجب أن تحقق بعض الشروط، ويمكننا وضع هذه الشروط بديهيات
لنظرية الجزء. فعلاقة الجزء، مثلًا، ترتيب جزئي على الكائنات: كل كائن
جزء من نفسه (وإن كان جزءًا \emph{غير فعلي})؛ ولا يمكن لكائنين مختلفين
أن يكون كل منهما جزءًا من الآخر؛ وجزء جزء كائن هو نفسه جزء من ذلك
الكائن. لاحظ أن «هو جزء من» تشبه بهذا المعنى «هو مجموعة جزئية من»،
لكنها لا تشبه «هو عنصر من»، إذ إن الأخيرة ليست انعكاسية ولا متعدية.
\begin{align*}
& \lforall[x][\Atom{\Obj P}{x,x}] \\
& \lforall[x][\lforall[y][((\Part{x}{y} \land \Part{y}{x})
      \lif \eq[x][y])]] \\
& \lforall[x][\lforall[y][\lforall[z][((\Part{x}{y} \land
        \Part{y}{z}) \lif \Part{x}{z})]]]\\
\intertext{وفوق ذلك، لكل كائنين مجموع في مبحث الجزء والكل (أي كائن
  يكون هذان الكائنان جزأين منه، ويكون أصغريًا من هذه الجهة).}  &
\lforall[x][\lforall[y][\lexists[z][\lforall[u][(\Part{z}{u} \liff
        (\Part{x}{u} \land \Part{y}{u}))]]]]
\end{align*}
وهذه ليست إلا بعض المبادئ الأساسية لعلاقة الجزء التي ينظر فيها
الميتافيزيقيون. غير أن صياغة مزيد من المبادئ أو كتابتها سرعان ما
تصعب من غير إدخال بعض العلاقات المعرّفة أولًا. فمعظم الميتافيزيقيين
المهتمين بمبحث الجزء والكل يعدون المبدأ الآتي صالحًا أيضًا: متى كان
للكائن~$x$ جزء فعلي~$y$، كان له كذلك جزء~$z$ لا يشترك مع~$y$ في أي
جزء، ويكون اندماج $y$ و$z$ هو~$x$.
\end{ex}

\end{document}
